%% ch-meta-k2.tex -- Part 2 (\ref{part:meta}), chapter 2: admissibility of the kernel rules over the core (K2).
%% Sources: structure-library/meta-*.scm, theorem-library/meta-*-laws.scm; the per-rule ledger
%% reference/K2-LEDGER.md.  Cites the design specification as \DS{n}.
\providecommand{\DS}[1]{\hyperref[ds:#1]{DS-#1}}
\providecommand{\VNBfig}[3]{\[\frac{\displaystyle #2}{\displaystyle #3}\;\;\mbox{\texttt{#1}}\]}
\providecommand{\VNBfigS}[3]{\[\frac{\displaystyle\begin{array}{c}#2\end{array}}{\displaystyle #3}\;\;\mbox{\texttt{#1}}\]}
\providecommand{\VNBnohyp}{\phantom{\Gamma}}
\providecommand{\FV}{\mathrm{FV}}
%% A core figure in a table cell: \Rl{hypotheses}{conclusion}.
\providecommand{\Rl}[2]{$\frac{#1}{#2}$}
\providecommand{\Sq}{\Rightarrow}

\chapter{Admissibility over the core}
\label{ch:meta-k2}

This chapter states what is proven, inside VNB, about the kernel's
rules.  The route is K2 (\DS{8}): formulas and sequents are encoded as
sets, a small system of rules called the \emph{core} is defined, and each
non-oracle kernel rule is proven \emph{admissible} over it.  What has to
be trusted is thereby reduced from the 57 checked relations to the core.

The definitions are in \texttt{structure-library/meta-*.scm}, the
theorems in the files of the same names in \texttt{theorem-library/},
all in the library \texttt{meta}.  Every theorem bills
\texttt{modulo 0}.

%% =========================================================
\section{The encoding}\label{sec:meta-encoding}

The rules are schemas, and VNB has no formula variables, so formulas
are VNB objects.  The encoding is locally nameless (\DS{10}): a bound
variable is a de~Bruijn index, a free variable a name.
Alpha-equivalence is then equality, and substitution for a free name is
structural.

Every node is a literal list whose first entry is a numeral tag:
\[
\begin{gathered}
\mathrm{EFVAR}(n)=[0,n],\quad \mathrm{EBVAR}(k)=[1,k],\quad
\mathrm{ECONST}(c)=[2,c],\\
\mathrm{EAPP}(h,l)=[3,h,l],\quad \mathrm{EBIND}(h,n,a,b)=[4,h,n,a,b].
\end{gathered}
\]
$h$ is a numeral head code, $l$ a list of nodes, $a$ the list of
arguments outside a binder (the domain of \texttt{sep},
\texttt{big-union}, \texttt{vnb-lambda}) and $b$ its body.

$\mathrm{EXPR}$ is a union over $\mathtt{nn}$ of stages.  The free
names $\FV$, the opening of an index $\mathrm{OPEN}$, the substitution
$\mathrm{SUBST}$ and local closure $\mathrm{LC}$ are defined by
recursion on the stage, with the usual substitution lemmas.  A
well-formedness predicate separates terms from formulas by head code.
$\mathrm{CLF}$ is the set of sentences (closed formulas), $\mathrm{CLT}$
the set of closed terms.

Predicate constants have codes ($\mathrm{PRED\text{-}CODE}$).  A
functoid application is $\mathrm{EAPP}(h,l)$ with $h\geq 32$, read
through a table $\mathrm{FUN\text{-}DEF}$ of rows $(h, n, b)$
(\DS{44}).  The tables hold the rows the META theorems need; filling
them for the whole library is the translator's task (K3).

In what follows formulas are written in ordinary notation for their
encodings: $A \wedge B$ for $\mathrm{ENC\text{-}AND}(A,B)$, $B[x:=t]$ for
$\mathrm{OPEN}(B,0,t)$, and so on.

%% =========================================================
\section{Theory, sequents and derivability}\label{sec:meta-derivable}

Fix a \emph{theory} $T$: a set of encoded sentences, standing for the
axioms of the catalog and the theorems installed at the time.  A cited
theorem is a member of $T$ and nothing more.

A \emph{sequent} $\Gamma \Longrightarrow A$ is a pair $[\Gamma, A]$ with
$\Gamma$ a finite set of sentences ($\mathrm{CTX\text{-}SET}$) and $A$ a
sentence; $\mathrm{SEQ\text{-}SET}$ is the set of sequents.  Contexts are
sets, so ``$\Gamma, A$'' is $\Gamma \cup \{A\}$.

Each core rule is a relation $\mathrm{CR\text{-}}x(T, \mathit{hs}, c)$
between a list $\mathit{hs}$ of hypothesis sequents and a conclusion
sequent $c$.  $\mathrm{CORE\text{-}RULE}(T, \mathit{hs}, c)$ is their
disjunction.  $\mathrm{DERIVABLE}(T)$ is the least subset of
$\mathrm{SEQ\text{-}SET}$ closed under
$\mathrm{CORE\text{-}RULE}(T,\cdot,\cdot)$.

We write $\vdash_T \Gamma \Longrightarrow A$ for
$[\Gamma, A] \in \mathrm{DERIVABLE}(T)$.  Each core rule has a closure
law, generated in \texttt{meta-core-laws.scm} and named
\texttt{derivable-}$x$: if the hypotheses of an instance are derivable,
so is its conclusion.

%% =========================================================
\section{The core}\label{sec:meta-core}

The core is chosen so that each member is a rule of natural deduction
with its side condition, or a one-line schema that a reader checks
against the intended meaning of a constructor (\DS{40}).  A relation
enters the core only by a design decision, with its claim written beside
it in \texttt{meta-core.scm} (\DS{14}).

$\mathrm{CORE\text{-}RULE}$ has 56 disjuncts.  They are listed below as
figures, hypotheses above the line.  A side condition follows in
brackets.  In the set-formation and definedness schemas, $p$ or $b$
carries the premise $\mathrm{NBG\text{-}OK}$ (every quantifier confined
to sets, read through the tables), written [NBG] (\DS{39}).

{\small
\begin{longtable}{@{}p{0.22\linewidth}p{\dimexpr0.78\linewidth-2\tabcolsep\relax}@{}}
\toprule
core rule & figure \\
\midrule
\endhead
\multicolumn{2}{@{}l}{\emph{Logic} (21)} \\
and-intro & \Rl{\Gamma\Sq A \quad \Gamma\Sq B}{\Gamma\Sq A\wedge B} \\
and-elim & \Rl{\Gamma, A, B\Sq G}{\Gamma\Sq G}\ \ [$A\wedge B\in\Gamma$] \\
or-intro-left, -right & \Rl{\Gamma\Sq A}{\Gamma\Sq A\vee B}\ \ \Rl{\Gamma\Sq B}{\Gamma\Sq A\vee B} \\
or-elim & \Rl{\Gamma,A\Sq G \quad \Gamma,B\Sq G}{\Gamma\Sq G}\ \ [$A\vee B\in\Gamma$] \\
implies-intro & \Rl{\Gamma,A\Sq B}{\Gamma\Sq A\to B} \\
backchain & \Rl{\Gamma\Sq A}{\Gamma\Sq G}\ \ [$A\to G\in\Gamma$] \\
not-intro & \Rl{\Gamma,A\Sq\textsc{falsity}}{\Gamma\Sq\neg A} \\
not-elim & \Rl{\VNBnohyp}{\Gamma\Sq G}\ \ [$\neg A\in\Gamma$, $A\in\Gamma$] \\
iff-intro & \Rl{\Gamma,A\Sq B \quad \Gamma,B\Sq A}{\Gamma\Sq A\leftrightarrow B} \\
iff-elim & \Rl{\Gamma, A\to B, B\to A\Sq G}{\Gamma\Sq G}\ \ [$A\leftrightarrow B\in\Gamma$] \\
forall-intro & \Rl{\Gamma\Sq B[x:=y]}{\Gamma\Sq\forall x.\,B}\ \ [$y\notin\FV(\Gamma)\cup\FV(B)$] \\
forall-elim & \Rl{\Gamma, B[x:=t]\Sq G \quad \Gamma\Sq t=t}{\Gamma\Sq G}\ \ [$\forall x.\,B\in\Gamma$, $t\in\mathrm{CLT}$] \\
forsome-intro & \Rl{\Gamma\Sq B[x:=t] \quad \Gamma\Sq t=t}{\Gamma\Sq\exists x.\,B}\ \ [$t\in\mathrm{CLT}$] \\
forsome-elim & \Rl{\Gamma, B[x:=y]\Sq G}{\Gamma\Sq G}\ \ [$\exists x.\,B\in\Gamma$, $y\notin\FV(\Gamma)\cup\FV(G)$] \\
assumption & \Rl{\VNBnohyp}{\Gamma\Sq A}\ \ [$A\in\Gamma$] \\
theorem-assumption & \Rl{\VNBnohyp}{\Gamma\Sq A}\ \ [$A\in T$] \\
cut & \Rl{\Gamma\Sq A \quad \Gamma,A\Sq G}{\Gamma\Sq G} \\
weakening & \Rl{\Delta\Sq G}{\Gamma\Sq G}\ \ [$\Delta\subseteq\Gamma$] \\
proof-by-contradiction & \Rl{\Gamma,\neg G\Sq\textsc{falsity}}{\Gamma\Sq G} \\
truth-intro & \Rl{\VNBnohyp}{\Gamma\Sq\textsc{truth}} \\
\midrule
\multicolumn{2}{@{}l}{\emph{Equality} (4)} \\
eq-subst, eq-subst-sym & \Rl{\Gamma\Sq P[z:=t]}{\Gamma\Sq P[z:=s]}\ \ [one of $s=t$, $s\equiv t$, $t=s$, $t\equiv s$ in $\Gamma$] \\
reflexivity & \Rl{\VNBnohyp}{\Gamma\Sq t=t}\ \ [$t$ a variable or an atomic constant] \\
quasi-reflexivity & \Rl{\VNBnohyp}{\Gamma\Sq t\equiv t} \\
\midrule
\multicolumn{2}{@{}l}{\emph{Structural membership and tuples} (7)} \\
union-mem & \Rl{\Gamma\Sq S_1=S_1 \ \cdots\ \Gamma\Sq S_n=S_n}{\Gamma\Sq x\in\bigcup(S_1..S_n)\leftrightarrow x\in S_1\vee\cdots\vee x\in S_n} \\
intersection-mem & \Rl{\VNBnohyp}{\Gamma\Sq x\in\bigcap(S_1..S_n)\leftrightarrow x\in S_1\wedge\cdots\wedge x\in S_n}\ \ [$n\geq 1$] \\
cartesian-mem & \Rl{\VNBnohyp}{\begin{array}{@{}l@{}}\Gamma\Sq x\in\mathrm{CART}(A_1..A_n)\leftrightarrow{}\\ \quad x\in\mathrm{TUPLES}(\mathrm{SET})\wedge\mathrm{LENGTH}(x)=n\wedge{}\\ \quad\bigwedge_i\mathrm{NTH}(i,x)\in A_i\end{array}}\ \ [$n\geq1$] \\
tuples-mem & \Rl{\Gamma\Sq A=A}{\Gamma\Sq[e_1..e_n]\in\mathrm{TUPLES}(A)\leftrightarrow\bigwedge_i e_i\in A} \\
nth-eq & \Rl{\Gamma\Sq L=L}{\Gamma\Sq\mathrm{NTH}(k,L)\equiv e_k}\ \ [$L=[e_1..e_n]$, $1\leq k\leq n$] \\
length-eq & \Rl{\Gamma\Sq L=L}{\Gamma\Sq\mathrm{LENGTH}(L)\equiv n} \\
tuple-eq & \Rl{\Gamma\Sq L=L \quad \Gamma\Sq M=M}{\Gamma\Sq L=M\leftrightarrow\bigwedge_i a_i=b_i}\ \ [$L=[a_1..a_n]$, $M=[b_1..b_n]$] \\
\midrule
\multicolumn{2}{@{}l}{\emph{Conditional terms} (2)} \\
if-true & \Rl{\Gamma\Sq p \quad \Gamma, \mathrm{IF}(p,a,b)\equiv a\Sq G}{\Gamma\Sq G} \\
if-false & \Rl{\Gamma\Sq \neg p \quad \Gamma, \mathrm{IF}(p,a,b)\equiv b\Sq G}{\Gamma\Sq G} \\
\midrule
\multicolumn{2}{@{}l}{\emph{Set formation, descriptions, functions} (9)} \\
sep-sethood & \Rl{\Gamma\Sq A\in\mathrm{SET}}{\Gamma\Sq\mathrm{SEP}\,x{\in}A.\,p\in\mathrm{SET}}\ \ [NBG] \\
sep-mem & \Rl{\VNBnohyp}{\Gamma\Sq t\in\mathrm{SEP}\,x{\in}A.\,p\leftrightarrow t\in A\wedge p[x:=t]}\ \ [NBG] \\
comp-mem & \Rl{\VNBnohyp}{\Gamma\Sq t\in\mathrm{COMP}\,x.\,p\leftrightarrow t\in\mathrm{SET}\wedge p[x:=t]}\ \ [NBG] \\
big-union-sethood & \Rl{\Gamma\Sq A\in\mathrm{SET} \quad \Gamma\Sq\forall z.\,(z\in A\to b\in\mathrm{SET})}{\Gamma\Sq\bigcup z{\in}A.\,b\in\mathrm{SET}}\ \ [NBG] \\
big-union-mem & \Rl{\VNBnohyp}{\Gamma\Sq x\in\bigcup z{\in}A.\,b\leftrightarrow\exists z.\,(z\in A\wedge x\in b)}\ \ [NBG] \\
iota-def & \Rl{\Gamma\Sq\exists x.\,(B\wedge\forall y.\,(B[x:=y]\to x=y)) \quad \Gamma, B[x:=\iota x.B]\Sq G}{\Gamma\Sq G} \\
iota-in-elim & \Rl{\Gamma\Sq\iota x.B=\iota x.B \quad \Gamma, B[x:=\iota x.B]\Sq G}{\Gamma\Sq G} \\
lambda-type & \Rl{\Gamma\Sq\forall x.\,(x\in A\to b\in D) \quad \Gamma\Sq A\in\mathrm{SET} \quad \Gamma\Sq D=D}{\Gamma\Sq\lambda x{\in}A.\,b\in\mathrm{FUN}(A,D)}\ \ [NBG; one binder] \\
lambda-beta-eq & \Rl{\Gamma\Sq t\in A \quad \Gamma\Sq b[x:=t]\in\mathrm{SET}}{\Gamma\Sq(\lambda x{\in}A.\,b)(t)\equiv b[x:=t]}\ \ [NBG] \\
\midrule
\multicolumn{2}{@{}l}{\emph{Definedness} (9; \DS{42})} \\
def-strict-in-l, -r & \Rl{\Gamma\Sq w\in A}{\Gamma\Sq w=w}\ \ \Rl{\Gamma\Sq w\in A}{\Gamma\Sq A=A} \\
def-total & \Rl{\Gamma\Sq t_1=t_1 \ \cdots\ \Gamma\Sq t_n=t_n}{\Gamma\Sq h(t_1..t_n)=h(t_1..t_n)}\ \ [$h$ a total constructor] \\
def-binder & \Rl{\Gamma\Sq A=A}{\Gamma\Sq Qx{\in}A.\,b=Qx{\in}A.\,b}\ \ [$Q$ one of SEP, BIG-UNION, VNB-LAMBDA; NBG] \\
def-comp & \Rl{\VNBnohyp}{\Gamma\Sq\mathrm{COMP}\,x.\,p=\mathrm{COMP}\,x.\,p}\ \ [NBG] \\
def-list & \Rl{\Gamma\Sq e_1\in\mathrm{SET}\ \cdots\ \Gamma\Sq e_n\in\mathrm{SET}}{\Gamma\Sq[e_1..e_n]=[e_1..e_n]} \\
def-strict-arg & \Rl{\Gamma\Sq h(t_1..t_n)=h(t_1..t_n)}{\Gamma\Sq t_i=t_i}\ \ [$h$ a strict constructor] \\
def-strict-dom & \Rl{\Gamma\Sq Qx{\in}A.\,b=Qx{\in}A.\,b}{\Gamma\Sq A=A} \\
def-strict-pred & \Rl{\Gamma\Sq P(t_1..t_n)}{\Gamma\Sq t_i=t_i}\ \ [$P$ a predicate constant] \\
\midrule
\multicolumn{2}{@{}l}{\emph{Congruence, functoids, tuples} (4; \DS{43}--\DS{45})} \\
iota-cr & \Rl{\Gamma\Sq B[x:=y]\leftrightarrow C[x:=y]}{\Gamma\Sq\iota x.B\equiv\iota x.C}\ \ [$y$ fresh] \\
tbind-cr & \Rl{\Gamma\Sq B[x:=y]\leftrightarrow C[x:=y]}{\Gamma\Sq Qx{\in}A.\,B\equiv Qx{\in}A.\,C}\ \ [$Q$ SEP or COMP; for BIG-UNION, VNB-LAMBDA the premise is $b[y]\equiv c[y]$; $y$ fresh] \\
functoid-beta-eq & \Rl{\VNBnohyp}{\Gamma\Sq F(l)\equiv b[l]}\ \ [$(F,|l|,b)$ a row of FUN-DEF] \\
tuple-eta & \Rl{\Gamma\Sq x\in\mathrm{CART}(A_1..A_n)}{\Gamma\Sq x=[\mathrm{NTH}(1,x)..\mathrm{NTH}(n,x)]}\ \ [$n\geq1$] \\
\bottomrule
\end{longtable}
}

Where a strict conclusion would assert that a term denotes, the core
concludes a quasi-equation instead: the conditional rules, lambda-beta,
the two congruence rules and functoid beta.  Where the kernel reads the
definedness of a term off its certificate, the core states it as a
premise: forall-elim, forsome-intro, iota-in-elim and the list schemas.
The differences between such a core rule and the kernel's figure are the
bridges of Section~\ref{sec:meta-bridges}.

A schema whose conclusion makes a class defined that its other premises
do not carries that class's definedness as a premise (\DS{53}):
tuples-mem the class $A$, union-mem every entry $S_i$, lambda-type the
codomain $D$.  $\mathrm{TUPLES}$, $\mathrm{UNION}$ and $\mathrm{FUN}$ are
strict in these arguments, and without the premise the schema is false
at the empty tuple, at a non-denoting entry, or over an empty domain.
lambda-beta-eq carries the value's sethood, $b[x:=t]\in\mathrm{SET}$: a
pair with a proper class as its second component does not denote.

Each schema beyond the logic, equality and structural groups carries its
claim beside it in \texttt{meta-core.scm}; where a schema is sound only
under a precondition, the precondition is written there and pinned by a
suite control.  The strictness of a predicate application, for example,
is sound only because every instantiation of a definition owes the
definedness of the instance (\DS{19}, \DS{42}).  The suite checks that
the four schemas above carry their premises.

%% =========================================================
\section{The statement of K2}\label{sec:meta-k2-statement}

For each non-oracle kernel rule $R$, $\mathrm{REL\text{-}}R(T,
\mathit{hs}, c)$ is the relation its checker enforces, transcribed by
hand into VNB.  Each transcription is kept beside the checker's line
numbers, to be reviewed as a pair (\DS{9}).

An instance of $R$ is a figure
\[
  \frac{\Gamma_1 \Longrightarrow A_1 \qquad \cdots \qquad \Gamma_n \Longrightarrow A_n}
       {\Gamma \Longrightarrow A}\;\;\mbox{\textit{R}}
\]
with $\mathrm{REL\text{-}}R(T, [\Gamma_1\Sq A_1, \ldots,
\Gamma_n\Sq A_n], \Gamma\Sq A)$.

\begin{dfn}
$R$ is \emph{admissible over the core} of $T$ when every instance with
derivable hypotheses has a derivable conclusion:
\[
  \frac{\vdash_T \Gamma_1 \Longrightarrow A_1 \qquad \cdots \qquad \vdash_T \Gamma_n \Longrightarrow A_n}
       {\vdash_T \Gamma \Longrightarrow A}
\]
\end{dfn}

\begin{claim}[K2, per rule]\label{claim:meta-k2}
Every non-oracle kernel rule is admissible over the core.  As a VNB
theorem, for a rule $R$:
\[
\begin{aligned}
&\forall T, \mathit{hs}, c.\;
\mathrm{REL\text{-}}R(T,\mathit{hs},c) \Rightarrow{}\\
&\qquad\bigl(\mathrm{HYPS\text{-}IN}(\mathit{hs}, \mathrm{DERIVABLE}(T)) \Rightarrow
c \in \mathrm{DERIVABLE}(T)\bigr).
\end{aligned}
\]
\end{claim}

$\mathrm{HYPS\text{-}IN}(\mathit{hs}, S)$ says that $\mathit{hs}$ is a
list, a set, every entry of which is in $S$ (\DS{47}).

For a rule that is itself a core rule, the claim is its closure law.
The content lies in the rules that are not: there the proof is a
derivation template, built by recursion on the syntax.  What the
theorems establish is that trusting the kernel's relations is no more
than trusting the core.  They do not establish that the core is sound
(link~3 of Section~\ref{sec:meta-chain}), nor that the Scheme of a
checker computes its relation (link~2).

\paragraph{Four kinds of proof.}
\begin{enumerate}
\item \emph{Derivation templates.}  A rule outside the core is derived
  by a template (\texttt{adm-detach}, \texttt{adm-union-intro},
  \texttt{adm-tuples-intro}, \ldots).
\item \emph{Contextual rewriting.}  The rules that rewrite at an
  arbitrary position (\texttt{macete}, \texttt{macete-hyp},
  \texttt{lambda-beta}, \texttt{lambda-beta-hyp}, \texttt{nth-reduce},
  \texttt{length-reduce}, \texttt{functoid-beta}, and the two
  decompositions) are reduced to a theorem CR: a rewrite valid in the
  local context of a position is valid at that position.  CR is proven by
  induction on positions, by stages: stages 1--5 cover the connectives,
  atoms, quantifiers, \texttt{if}, the arguments of a term and the domain
  of a binder; stage~6 the bodies of \texttt{sep}, \texttt{big-union} and
  \texttt{vnb-lambda}, through tbind-cr; stage~7 the bodies of
  \texttt{comp} and \texttt{iota}.  A rewrite at several positions is a
  chain of one-position rewrites, and its admissibility follows from CR
  by induction on the length of the chain
  (Section~\ref{sec:meta-positions}).
\item \emph{Definedness rows.}  Each clause of the certificate
  \texttt{pi--defined?} is transcribed as a relation
  $\mathrm{DEF\text{-}}x$ and derived over the core or admitted as a
  schema (Section~\ref{sec:meta-lemma-star}).
\item \emph{Inductions.}  The three induction rules are derived from the
  class-form induction axioms of the catalog and comp-mem, for a
  predicative body.
\end{enumerate}

%% =========================================================
\section{Coverage}\label{sec:meta-coverage}

Each of the 57 non-oracle rules has one of the following forms.

\begin{center}
\small
\begin{tabular}{@{}lp{\dimexpr0.88\linewidth-2\tabcolsep\relax}@{}}
\toprule
form & meaning \\
\midrule
C & The kernel figure is a core figure.  The theorem is the closure law. \\
B & The kernel figure differs from a core figure by a small step, the
  \emph{bridge}: the kernel replaces a formula the core keeps, or omits a
  premise the core states.  The kernel relation is transcribed and its
  admissibility proven (\texttt{adm-}$R$\texttt{-k}). \\
B* & As B, when the premise the kernel omits is the definedness
  $\Gamma \Longrightarrow t = t$ of a term the kernel certifies.  The
  admissibility is proven for a term certified by a member of the context
  (\texttt{adm-}$R$\texttt{-k}) and for a term certified by any
  composition of the definedness rows of lemma~(*)
  (\texttt{adm-}$R$\texttt{-kd}; Section~\ref{sec:meta-lemma-star}).
  A certificate that uses an open row is not covered. \\
D & Not a core rule.  Transcribed and derived in full over the core
  (\texttt{adm-}$R$). \\
PS & A rewriting rule, admissible for a rewrite at any finite set of
  positions, at any stage of CR (\texttt{adm-}$R$\texttt{-s};
  Section~\ref{sec:meta-positions}). \\
\bottomrule
\end{tabular}
\end{center}

\begin{center}
\small
\begin{tabular}{@{}lrp{\dimexpr0.82\linewidth-4\tabcolsep\relax}@{}}
\toprule
form & count & rules \\
\midrule
C & 19 & and-intro, or-intro-left, or-intro-right, implies-intro,
  not-intro, iff-intro, forall-intro, not-elim, assumption,
  theorem-assumption, cut, weakening, backchain, proof-by-contradiction,
  eq-subst, quasi-reflexivity, sep-sethood, big-union-sethood, iota-def \\
B & 5 & and-elim, or-elim, iff-elim, forsome-elim, iota-in-elim \\
B* & 5 & forsome-intro, forall-elim, reflexivity, if-true, if-false \\
D & 19 & detach, cartesian-intro, cartesian-elim, tuples-intro,
  tuples-elim, union-intro, union-elim, intersection-intro,
  intersection-elim, nn-induction, transfinite-induction,
  transfinite-induction-3cases, sep-mem-intro, sep-mem-elim,
  comp-mem-intro, comp-mem-elim, big-union-mem-intro, big-union-mem-elim,
  lambda-type \\
PS & 9 & nth-reduce, length-reduce, functoid-beta, lambda-beta,
  lambda-beta-hyp, macete, macete-hyp, cartesian-decompose,
  tuple-equality-decompose \\
\midrule
all & 57 & \\
\bottomrule
\end{tabular}
\end{center}

So Claim~\ref{claim:meta-k2} is proven in full for the C, B, D and PS
rows, and for the B* rows for every certificate built from the rows of
lemma~(*).  One qualification applies to the D rows tuples-intro,
union-intro and lambda-type, which owe the definedness of a class
(Section~\ref{sec:meta-owed}).  The D theorem is for the figure with
the owed premise.  Where the kernel omits the premise because its
certificate accepts the class, the figure is covered at the empty tuple
and for lambda-type at one binder, for a class certified by a member of
the context; it is open for union-intro and for lambda-type at a binder
list.

The ledger \texttt{reference/K2-LEDGER.md} lists, for each of the 57
rules, its form, the relation, the theorems that discharge it and what
is open.

%% =========================================================
\section{The bridged rules}\label{sec:meta-bridges}

For each B and B* rule: the kernel figure, the core figure, the bridge,
and the theorem.  The kernel relations are transcribed in
\texttt{structure-library/meta-nbg.scm}, the theorems proven in section~B
of \texttt{theorem-library/meta-nbg-laws.scm}.  $\Gamma'$ is the context
with the eliminated formula removed.

\subsection{The kernel replaces, the core keeps}

The kernel's elimination rules remove the formula they eliminate from
the context; the core's keep it.  The bridge is one weakening.

\VNBfig{and-elim (kernel)}{\Gamma', A, B \Longrightarrow G}{\Gamma \Longrightarrow G}
\VNBfig{and-elim (core)}{\Gamma, A, B \Longrightarrow G}{\Gamma \Longrightarrow G}
with $A \wedge B \in \Gamma$ and $\Gamma' = \Gamma \setminus \{A\wedge B\}$.
Bridge: $\Gamma', A, B \subseteq \Gamma, A, B$, so weakening takes the
kernel's hypothesis to the core's.  Theorem \texttt{adm-and-elim-k}.

\VNBfig{or-elim (kernel)}{\Gamma', A \Longrightarrow G \qquad \Gamma', B \Longrightarrow G}{\Gamma \Longrightarrow G}
with $A \vee B \in \Gamma$.  Core: the same with $\Gamma$ for $\Gamma'$.
Bridge: weakening in each branch.  Theorem \texttt{adm-or-elim-k}.

\VNBfig{iff-elim (kernel)}{\Gamma', A \to B, B \to A \Longrightarrow G}{\Gamma \Longrightarrow G}
with $A \leftrightarrow B \in \Gamma$.  Bridge: weakening.  Theorem
\texttt{adm-iff-elim-k}.

\VNBfig{forsome-elim (kernel)}{\Gamma', B[x:=y] \Longrightarrow G}{\Gamma \Longrightarrow G}
with $\exists x.\,B \in \Gamma$ and $y \notin \FV(\Gamma') \cup \FV(G)
\cup \FV(\exists x.\,B)$.  The core asks $y \notin \FV(\Gamma) \cup
\FV(G)$.  Bridge: a free name of a member of $\Gamma$ is free in a member
of $\Gamma'$ or in $\exists x.\,B$ (\texttt{ctx-fv-minus}), so the
kernel's eigenvariable condition gives the core's; then weakening.
Theorem \texttt{adm-forsome-elim-k}.

\subsection{The kernel reads a definedness off the context}

\VNBfig{iota-in-elim (kernel)}{\Gamma, B[x:=\iota x.B] \Longrightarrow G}{\Gamma \Longrightarrow G}
with $\iota x.B \in A$, $\iota x.B = u$ or $u = \iota x.B$ in $\Gamma$.
\VNBfig{iota-in-elim (core)}{\Gamma \Longrightarrow \iota x.B = \iota x.B \qquad \Gamma, B[x:=\iota x.B] \Longrightarrow G}{\Gamma \Longrightarrow G}
Bridge: the core's first hypothesis is derived from the context member,
by assumption and the strictness of $\in$ (def-strict-in-l) or of $=$
(\texttt{adm-def-strict-eq-l}, \texttt{-r}).  Theorem
\texttt{adm-iota-in-elim-k}.

\subsection{The kernel omits a definedness premise it certifies}

For the five B* rules the kernel omits the premise
$\Gamma \Longrightarrow t = t$ when its certificate accepts $t$ in
$\Gamma$.  When the certificate declines, the kernel posts that premise
and its figure is the core's: that case is C.

Each B* rule has two transcribed kernel forms.  In the form
$\mathrm{REL\text{-}}R\mathrm{\text{-}K}$ the term is certified by one
member of the context; in the form $\mathrm{REL\text{-}}R\mathrm{\text{-}KD}$
the sequent $\Gamma \Longrightarrow t = t$ is a member of
$\mathrm{DEF\text{-}DERIVABLE}(T)$, the sequents the definedness rows
reach (Section~\ref{sec:meta-lemma-star}).  The theorems are
\texttt{adm-}$R$\texttt{-k} and \texttt{adm-}$R$\texttt{-kd}.

\VNBfig{forall-elim (kernel, $t$ certified)}{\Gamma, B[x:=t] \Longrightarrow G}{\Gamma \Longrightarrow G}
\VNBfig{forall-elim (core)}{\Gamma, B[x:=t] \Longrightarrow G \qquad \Gamma \Longrightarrow t = t}{\Gamma \Longrightarrow G}
with $\forall x.\,B \in \Gamma$.  Bridge: $\vdash_T \Gamma \Longrightarrow
t=t$ whenever $\Gamma$ certifies $t$; this is lemma~(*).  Theorems
\texttt{adm-forall-elim-k} and \texttt{adm-forall-elim-kd}.

\VNBfig{forsome-intro (kernel, $t$ certified)}{\Gamma \Longrightarrow B[x:=t]}{\Gamma \Longrightarrow \exists x.\,B}
\VNBfig{forsome-intro (core)}{\Gamma \Longrightarrow B[x:=t] \qquad \Gamma \Longrightarrow t = t}{\Gamma \Longrightarrow \exists x.\,B}
The kernel accepts $t$ certified by the instance itself (a strict atom
among the sentences $B[x:=t]$ makes true) or by the context.  Bridge:
from the instance, by cut with the walks \texttt{adm-true-walk} and
\texttt{adm-strict-atom}; from the context or the rows, lemma~(*).  Theorems
\texttt{adm-forsome-intro-k} and \texttt{adm-forsome-intro-kd}.

\VNBfig{reflexivity (kernel)}{\VNBnohyp}{\Gamma \Longrightarrow t = t}
with $\Gamma$ certifying $t$.  The core's reflexivity fires only for a
variable or an atomic constant.  Bridge: lemma~(*).  Theorems
\texttt{adm-reflexivity-k} and \texttt{adm-reflexivity-kd}.  The
reflexivity checker also accepts $\Gamma \Longrightarrow \textsc{truth}$
with no hypothesis; that figure is the core's truth-intro (theorem
\texttt{adm-reflexivity-kt}).  The kernel has no other operation that
concludes \textsc{truth}.

\VNBfig{if-true (kernel, $a$ certified)}{\Gamma \Longrightarrow p \qquad \Gamma, \mathrm{IF}(p,a,b) = a \Longrightarrow G}{\Gamma \Longrightarrow G}
\VNBfig{if-true (core)}{\Gamma \Longrightarrow p \qquad \Gamma, \mathrm{IF}(p,a,b) \equiv a \Longrightarrow G}{\Gamma \Longrightarrow G}
Bridge: from $\Gamma \Longrightarrow a = a$ (lemma~(*)) and the
quasi-equation the core hands over, the strict equation (a
strictness law of $=$), then a cut.  The kernel hands over the strict
equation only when $a$ is certified, otherwise the quasi-equation
(\DS{20}).  Theorems \texttt{adm-if-true-k} and
\texttt{adm-if-true-kd}; if-false likewise with $\neg p$ and $b$,
\texttt{adm-if-false-k} and \texttt{adm-if-false-kd}.

\subsection{A hypothesis that adds nothing}

The checkers of forall-elim, forsome-elim, if-true, if-false and
iota-in-elim also accept a hypothesis $\Delta \Longrightarrow G$ with the
conclusion's goal and $\Delta \subseteq \Gamma$: the formula the rule
would add is already in the context.  That figure is one weakening
(theorem \texttt{adm-adds-nothing-k}).

When the bound variable does not occur in the body, forsome-elim imposes
no eigenvariable condition.  That instance is an instance of the kernel
form for any fresh $y$, since opening a vacuous body does not change it;
the vacuous forall-elim is the kernel form at a fresh variable $y$, which
the rows certify by reflexivity.

%% =========================================================
\section{Rules that owe the definedness of a class}\label{sec:meta-owed}

Three kernel rules conclude a membership, or a $\mathrm{FUN}$ statement,
that makes a class defined which their other hypotheses do not
(\DS{53}).  Each owes the definedness of that class unless its
certificate accepts the class in $\Gamma$.
\begin{itemize}
\item \texttt{tuples-intro} at the empty tuple owes
  $\Gamma \Longrightarrow A = A$.  With $n \geq 1$ components the
  hypothesis $\Gamma \Longrightarrow a_1 \in A$ makes $A$ defined, and
  nothing is owed.
\item \texttt{union-intro} $k$ owes $\Gamma \Longrightarrow S_j = S_j$ for
  each entry $S_j$ other than $S_k$; the hypothesis
  $\Gamma \Longrightarrow x \in S_k$ makes $S_k$ defined.
\item \texttt{lambda-type} owes $\Gamma \Longrightarrow B = B$ for the
  codomain $B$.  Over an empty domain the typing hypothesis holds for
  every $B$.
\end{itemize}
The owed sequents are hypotheses of the inference, and the checker
accepts each omitted one only when \texttt{chk-elim-certified?} accepts
the class.  The core schemas tuples-mem, union-mem and lambda-type carry
the same premises (Section~\ref{sec:meta-core}).

\paragraph{The relations.}
$\mathrm{REL\text{-}TUPLES\text{-}INTRO}$ is the figure with $n \geq 1$
components, and $\mathrm{REL\text{-}TUPLES\text{-}INTRO\text{-}0}$ the
empty tuple with the owed premise (\texttt{meta-core.scm}).
$\mathrm{REL\text{-}UNION\text{-}INTRO}$ is the owed form: the
membership, then one definedness sequent for each entry.
$\mathrm{REL\text{-}LAMBDA\text{-}TYPE\text{-}1}$ to
$\text{-}4$ and $\mathrm{REL\text{-}LAMBDA\text{-}TYPE\text{-}N}$
(\texttt{meta-ind.scm}) carry $\Gamma \Longrightarrow B = B$ as their last
hypothesis.  The certified forms
$\mathrm{REL\text{-}TUPLES\text{-}INTRO\text{-}0K}$ and
$\mathrm{REL\text{-}LAMBDA\text{-}TYPE\text{-}K}$ (one binder;
\texttt{meta-nbg.scm}) drop the premise for a class certified by a
member of the context, as the -K forms of the B* rules do.

\paragraph{The theorems.}  \texttt{adm-tuples-intro},
\texttt{adm-tuples-intro-0} and \texttt{adm-tuples-intro-0k};
\texttt{adm-union-intro}; \texttt{adm-lambda-type-1} to \texttt{-4},
\texttt{adm-lambda-type-n} and \texttt{adm-lambda-type-k}.  The closure
law used by \texttt{adm-union-intro}, \texttt{derivable-union-mem-c},
assumes the definedness of every entry derivable.
\texttt{derivable-union-mem-u} obtains those premises from the
definedness of the union itself, by the strictness of $\mathrm{UNION}$
in its arguments.  The decomposition lemmas \texttt{der-union-decompose}
and \texttt{der-union-binary} carry the same assumption.

\paragraph{The rewriting rules.}  The local instances of the union
decomposition, $\mathrm{J\text{-}UDEC}$, require each entry's definedness
derivable in the local context.  Those of lambda-beta,
$\mathrm{J\text{-}LB}$, require the value's sethood derivable there: the
kernel's \texttt{lambda-beta} owes $\Gamma \Longrightarrow b[x:=u] \in
\mathrm{SET}$ when the value may be a proper class, and the core's
lambda-beta-eq has it as a premise.

\paragraph{What is open.}  The kernel omits an owed sequent that its
certificate accepts.  For \texttt{union-intro} at any entry, and for
\texttt{lambda-type} at a binder list, that omission is not
transcribed.  The -K forms cover a class certified by one member of the
context; no form covers a class certified through a composition of the
rows of lemma~(*).

\subsection{lambda-type at a binder list}

The core has one-binder lambdas only.  A kernel term
$\mathrm{vnb\text{-}lambda}([x_1..x_n], A_1 \times \cdots \times A_n, b)$
is encoded as a one-binder lambda over the product,
$\mathrm{ENC\text{-}LAM}(\mathrm{ENC\text{-}CART}(a), c)$, where $c$ is $b$
with each $x_i$ replaced by the projection $\mathrm{NTH}(i, \cdot)$ of the
bound variable.  The relation fixes $c$ by the one property the rule
needs: opened at a name $y$, $c$ is $b$ opened at the projections of $y$,
in the order in which instantiating the $n$ guarded universals opens
them.

\texttt{adm-lambda-type-n} proves the kernel figure admissible at every
length $n \geq 1$; the fixed lengths 1 to 4 are proven separately in
literal form.  For a fresh $z$ in the product, the chain of guarded
universals is instantiated at the projections of $z$ by induction on the
number of guards left (\texttt{ltch-inst}).  This gives $c$ opened at $z$
in $B$.  Then implies-intro, forall-intro at $z$, and the core
lambda-type.

The theorem is relative to the encoding of a destructuring lambda,
which is the translator's (K3).  Two of its properties are argued in
the file, not proven: every member of a product is the list of its
projections (the core's tuple-eta says it), and $c$ is predicative
exactly when $b$ is.

%% =========================================================
\section{Rewriting at several positions}\label{sec:meta-positions}

The kernel's rewriting rules rewrite at every position their walk
explains, in one inference.  The META relations read such an inference
as a chain of one-position rewrites.
$\mathrm{REWRS}(J)$ is the least set of rewrite tuples that contains
each one-position rewrite justified by $J$ (at any stage of CR) and the
empty rewrite, and is closed under composition.
$\mathrm{REL\text{-}}R\mathrm{\text{-}S}$ is the rule $R$ over the $J$
set of its one-position form, rewritten by a member of
$\mathrm{REWRS}(J)$ (\texttt{meta-nbg.scm}).

The chain is the kernel's relation when the positions are taken in the
order of the kernel's walk.  The macete walk stops at the first
explained position, so the positions are disjoint.  The local contexts
are read off different sides at a connective: at $A \wedge B$ the right
conjunct is rewritten under the original left, the left under the
rewritten right; at $A \vee B$ the left disjunct under the negation of
the original right, the right under the negation of the rewritten left;
at $A \to B$ the consequent under the original antecedent.  Taking the
right conjunct, the left disjunct and the consequent first, the context
the $k$-th step computes is the walk's.  The reduction rules read no
local context: their reachability, created redexes and contractions
inside a contractum included, is a chain of single contractions.

\texttt{adm-rewrite-s} and \texttt{adm-rewrite-hyp-s} prove, by the
least-set induction of $\mathrm{REWRS}$ from the one-position law, that a
rewrite at any finite set of positions is admissible for a goal and for
an assumption.  The rows are \texttt{adm-}$R$\texttt{-s} for nth-reduce,
length-reduce, functoid-beta, lambda-beta and its hypothesis form,
macete and macete-hyp (with their several-variable forms), the union and
intersection decompositions, cartesian-decompose and
tuple-equality-decompose.

Composition adds no restriction, because each step's justification is
judged on the formula the earlier steps produced.  What stays outside
is what the one-position forms exclude: the licence of lambda-beta
through its binder's own scope, the binder-domain reads of the list
definedness checks, and any justification outside the rule's $J$ set.

%% =========================================================
\section{The definedness lemma (*)}\label{sec:meta-lemma-star}

\begin{lem}[(*)]
If the certificate \texttt{pi--defined?} accepts the term $t$ in the
context $\Gamma$, and every clause it uses is a row of
$\mathrm{DEF\text{-}RULE}$, then $\vdash_T \Gamma \Longrightarrow t = t$.
\end{lem}

The lemma is the bridge of every B* rule.  Each clause of the
certificate is transcribed as a relation $\mathrm{DEF\text{-}}x$
(\texttt{structure-library/meta-def.scm}), and each relation is either a
core schema (its closure law) or derived over the core.  A clause that
certifies a compound term from certified parts is the corresponding
relation applied to the parts.

$\mathrm{DEF\text{-}RULE}(T, \mathit{hs}, c)$ is the disjunction of the
derived and admitted rows of the table below, together with the three
core rules that start a reading: assumption (a member of the context),
theorem-assumption (a member of $T$, such as a definition a predicate
row reads) and reflexivity (an atom).  The rows $\mathrm{DEF\text{-}APP}$
and $\mathrm{DEF\text{-}LEN\text{-}TUPLES}$ enter only when $T$ is a set of sentences
holding the encoded axiom the row rests on.  $\mathrm{DEF\text{-}DERIVABLE}(T)$ is the
least set of sequents closed under $\mathrm{DEF\text{-}RULE}(T,\cdot,\cdot)$,
defined as $\mathrm{DERIVABLE}(T)$ is.

\begin{thm}[\texttt{def-derivable-sub}]
For every $T$,
\[
\mathrm{DEF\text{-}DERIVABLE}(T) \subseteq \mathrm{DERIVABLE}(T).
\]
\end{thm}

The proof is \texttt{def-closed-derivable}: $\mathrm{DERIVABLE}(T)$ is
closed under every row of $\mathrm{DEF\text{-}RULE}$, each row citing its
closure law.  The theorem is lemma~(*) for every composition of the
rows.  That the kernel's acceptance of $t$, through clauses that are
rows, puts $\Gamma \Longrightarrow t = t$ in
$\mathrm{DEF\text{-}DERIVABLE}(T)$ is a faithfulness claim, part of K3
(\DS{9}).

{\small
\begin{longtable}{@{}p{\dimexpr0.4\linewidth-2\tabcolsep\relax}p{\dimexpr0.28\linewidth-2\tabcolsep\relax}p{0.32\linewidth}@{}}
\toprule
clause & relation & status \\
\midrule
\endhead
atom (variable, constant) & reflexivity & core \\
$t \in \_$ or $\_ \in t$ in context & def-strict-in-l, -r & core \\
$t = \_$, $\_ = t$ in context & DEF-STRICT-EQ-L, -R & derived \\
$t$ outside binders in a true $\in$, $=$, $\leq$, $<$, or in a
  separation condition & DEF-STRICT-HYP & derived \\
structure hypothesis $\mathrm{IS\text{-}X}(t)$, $t \in X$ & DEF-STRUCT-HYP & derived \\
predicate hypothesis (any arity, leading existentials, nested, parent
  chain); accessor under $\mathrm{IS\text{-}X}(s)$ & DEF-PRED-HYP-N, DEF-CERT-HYP & derived \\
literal list of sets & def-list & core \\
total class constructor on certified arguments & def-total & core \\
\texttt{sep}, \texttt{big-union}, \texttt{vnb-lambda} over a certified
  domain, predicative body & def-binder & core \\
\texttt{comp}, predicative condition & def-comp & core \\
$\mathrm{NTH}(k, L)$, $L$ certified & DEF-NTH-LIST & derived \\
functoid with a certified instance; NTH or SEP through a functoid
  & DEF-FUN-INST, -NTH, -SEP, -STRICT & derived \\
$f(a)$ with $f \in \mathrm{FUN}(D,\_)$, $a \in D$ & DEF-APP & derived, $T$ holding \texttt{fun-domain-apply-def} \\
$\mathrm{LENGTH}(t)$ with $t \in \mathrm{TUPLES}(A)$ in context & DEF-LEN-TUPLES & derived, $T$ holding \texttt{length-in-nn} \\
$\mathrm{LENGTH}(t)$, $t$ typed in a class of tuples through a
  \texttt{sep} domain or a functoid body & -- & open \\
ground arithmetic; arithmetic on certified numbers & -- & open \\
applied structure operation on typed arguments & -- & open \\
view clause $\mathit{ACC}(V(r))$, the whole list certified (\DS{28}) & -- & no row \\
\bottomrule
\end{longtable}
}

The derived rows are proven in the laws files of
\texttt{meta-def.scm} and \texttt{meta-ind.scm}; the theorem for each is
in the ledger.  The open rows have no theorem.  Ground arithmetic needs
head codes for the arithmetic heads and their closure axioms in $T$; the
applied structure operation needs an encoding of an application to
several arguments and the slot's domain term; \texttt{LENGTH} through a
\texttt{sep} domain or a functoid body needs the certificate's walk to a
class of tuples transcribed.  All three are part of K3.  The view clause
is not admissible as a row.  The retired clause for $\mathrm{NTH}$ of a
literal list is not a row.  A certificate that uses an open clause is not
covered by the B* theorems (Chapter~\ref{ch:meta-limits}).

%% =========================================================
\section{Regression}\label{sec:meta-regression}

The META theorems are checked as every theorem is: by the certified load
by day and by the exam at night.  Their certificates follow the
definitions of the statement transitively (\DS{17}), so a change to a
core relation re-proves every META theorem that rests on it.

Each kernel repair comes with a reproducer and suite controls, at least
one of which fails on the code before the repair.  Each core schema with
a precondition has a control for it; each congruence and beta rule has a
control that it concludes only the quasi-equation.
